The cases \(B=0\) and \(d=0\) are immediate. If \(B=0\), every coefficient vanishes. If \(d=0\), every neighborhood is empty, every causal contrast is zero, and the score and both SNIPE estimators are zero under \({\rm def:zero-degree-conventions}\). Hence assume \(B>0\) and \(1\le d\le n\).

For the upper bound, set
\[
W_i:=Y_i^{\mathrm{obs}}g_{i,\beta,p}(Z),\qquad q_i:=|N_i|.
\]
After relabeling the coordinates in \(N_i\), the representer calculation in \({\rm lem:block-energy-representer}\), applied in dimension \(q_i\), gives
\[
\mathbb E_Z[W_i]=Y_i(\mathbf1)-Y_i(\mathbf0).
\]
Thus \(\widehat\tau_n^{\mathrm{SNIPE}}\) is design-unbiased for \(\tau_n\). The coefficient-mass condition gives \(|Y_i(z)|\le B\) for every \(i,z\). Centered-monomial orthogonality gives the exact local score energy
\[
\mathbb E_Z[g_{i,\beta,p}(Z)^2]
=A_i:=\sum_{r=1}^{\min\{\beta,q_i\}}\binom{q_i}{r}\frac{\Delta_r(p)^2}{[p(1-p)]^r}\le A_d.
\]
Every centered function of the binary vector \(Z_{N_i}\) has a unique Bernoulli Fourier expansion
\[
W_i-\mathbb E_ZW_i=\sum_{\varnothing\ne S\subseteq N_i}a_{i,S}\prod_{j\in S}(Z_j-p).
\]
By orthogonality and boundedness of the outcome,
\[
\sum_{\varnothing\ne S\subseteq N_i}a_{i,S}^2[p(1-p)]^{|S|}
=\operatorname{Var}_Z(W_i)
\le\mathbb E_Z[W_i^2]
\le B^2A_i.
\]
For nonempty \(S\), let \(I_S:=\{i:S\subseteq N_i\}\). Choosing any \(j\in S\), the out-degree condition gives \(|I_S|\le\#\{i:j\in N_i\}\le d\). This is the single-out-degree charge isolated by \({\rm lem:overlap-count}\). A second use of orthogonality followed by Cauchy--Schwarz yields
\[
\begin{aligned}
\operatorname{Var}_Z\!\left(\sum_{i\in V_n}W_i\right)
&=\sum_{S\ne\varnothing}[p(1-p)]^{|S|}\left(\sum_{i\in I_S}a_{i,S}\right)^2\\
&\le d\sum_{i\in V_n}\sum_{\varnothing\ne S\subseteq N_i}a_{i,S}^2[p(1-p)]^{|S|}\\
&\le ndB^2A_d.
\end{aligned}
\]
Consequently
\[
\mathbb E_Z[(\widehat\tau_n^{\mathrm{SNIPE}}-\tau_n)^2]
\le B^2\frac{dA_d}{n}
\le C_{\beta,p}^{(1)}B^2\frac{d\binom d{k_\star}}{n},
\]
where the last step is \({\rm lem:block-energy-representer}\). Also \(|\tau_n|\le B\). Projection onto \([-B,B]\) cannot increase distance to \(\tau_n\), and its squared distance is at most \(4B^2\). Therefore
\[
\sup_{(G_n,c)\in\mathcal M_{n,d,\beta}(B)}
\mathbb E_Z[(\widehat\tau_n^{\mathrm{up}}-\tau_n)^2]
\le C_{\beta,p}^{(2)}B^2\min\!\left\{1,\frac{d\binom d{k_\star}}{n}\right\}.
\]
The minimax upper inequality follows by using this single data-dependent estimator. In the nonsaturated regime, the preceding unclipped variance display and the minimax lower bound proved below show that unclipped SNIPE has the same rate.

For the lower bound, use the corrected block construction and abbreviate
\[
\kappa:=\min\{(2H_{\beta,p})^{-1},(4\pi)^{-1}\}.
\]
By \({\rm lem:block-family-admissibility}\), both priors are supported on the announced model class. Every active unit in block \(b\) has
\[
Y_i(z)=U_b+\sigma\delta h_d(z_{V_b}).
\]
Since \(L_d(h_d)=1\), its all-one versus all-zero contrast is \(\sigma\delta\). Inactive units contribute zero, so every schedule under \(\Pi_\sigma\) has
\[
\tau_n=\sigma\rho\delta.
\]
Thus the TTE separation is exactly \(2\rho\delta\).

Let \(\psi_s:=\sqrt{f_s}\), extended by zero outside \([-s,s]\). Since the cosine vanishes at both endpoints, \(\psi_s\) has a square-integrable weak derivative, and direct integration gives
\[
\|\psi_s'\|_2^2
=\frac{\pi^2}{4s^3}\int_{-s}^{s}\sin^2\!\left(\frac{\pi u}{2s}\right)du
=\frac{\pi^2}{4s^2}.
\]
For any translations \(a,b\), the fundamental theorem of calculus and Minkowski's inequality give
\[
\|\psi_s(\,\cdot-a)-\psi_s(\,\cdot-b)\|_2
\le |a-b|\,\|\psi_s'\|_2.
\]
Conditional on a block assignment \(Z_b=z\), the repeated observed block outcome is \(U_b+\sigma\delta h_d(z)\). Hence the conditional laws under the two signs are translates of \(f_s\) separated by \(2\delta h_d(z)\). With the convention \(H^2(P,Q)=\int(\sqrt{dP}-\sqrt{dQ})^2\), averaging over the common assignment law and using \(s=B/2\) and \(\mathbb E_Z[h_d(Z)^2]=A_d^{-1}\) yields
\[
H^2(P_{+,b},P_{-,b})
\le\frac{4\pi^2\delta^2}{B^2A_d}.
\]
Hellinger affinity tensorizes over the \(m\) independent active blocks, and \(1-\prod_b a_b\le\sum_b(1-a_b)\) for \(a_b\in[0,1]\). The inactive observations are a common factor. Therefore
\[
H^2(P_+,P_-)
\le\frac{4\pi^2m\delta^2}{B^2A_d}
\le4\pi^2\kappa^2\le\frac14.
\]
The last inequality follows from \(\kappa\le(4\pi)^{-1}\). Cauchy--Schwarz gives \(\|P_+-P_-\|_{\mathrm{TV}}\le H(P_+,P_-)\le1/2\).

For an arbitrary estimator, threshold at zero to test the two signs. With \(\theta:=\rho\delta\), a testing error forces squared estimation error at least \(\theta^2\). Hence
\[
\begin{aligned}
\max_{\sigma\in\{-1,1\}}
\mathbb E_{P_\sigma}[(\widehat\tau_n-\sigma\theta)^2]
&\ge\frac{\theta^2}{2}\{1-\|P_+-P_-\|_{\mathrm{TV}}\}\\
&\ge\frac{\rho^2\delta^2}{4}.
\end{aligned}
\]
Because the priors are supported on \(\mathcal M_{n,d,\beta}(B)\), the left-hand side is bounded above by the worst-case risk entering \(R_n^\star\). Now \(m=\lfloor n/d\rfloor\), so
\[
\frac12\le\rho=\frac{md}{n}\le1,
\qquad
\frac{A_d}{m}=\frac{dA_d/n}{\rho}.
\]
It follows that
\[
R_n^\star(d,\beta,p,B)
\ge\frac{\kappa^2}{16}B^2\min\!\left\{1,\frac{dA_d}{n}\right\}.
\]
Finally, if
\[
a_{\beta,p}:=\min_{\substack{1\le r\le\beta\\\Delta_r(p)\ne0}}
\frac{\Delta_r(p)^2}{[p(1-p)]^r}>0,
\]
then the order-\(k_\star\) summand gives \(A_d\ge a_{\beta,p}\binom d{k_\star}\). Thus one may take
\[
c_{\beta,p}:=\frac{\kappa^2}{16}\min\{1,a_{\beta,p}\}>0
\]
in the displayed lower bound, and enlarge the upper constant if necessary so that \(c_{\beta,p}\le C_{\beta,p}\). This proves all three risk inequalities, the exact block-mixture separation and Hellinger bound, and the asserted relation \(A_d/m=(dA_d/n)/\rho\), without requiring \(d\mid n\).